\documentclass[10pt,a4paper]{amsart}
\usepackage[margin=19mm]{geometry}
\usepackage{amsmath,amssymb,mathtools}
\usepackage[scr=rsfs,cal=euler]{mathalfa}
\usepackage{xfrac}
\usepackage[unicode,hidelinks,bookmarks=false]{hyperref}
\newcommand{\ie}{i.e.\ }
\newcommand{\cf}{cf.\ }
\newcommand{\resp}{respectively\ }
\newcommand{\Subsection}{\subsection}
\providecommand{\nobreakdash}{\nobreak\hbox{-}}
\setlength{\emergencystretch}{2em}
\multlinegap0pt
\usepackage{amssymb}
\usepackage{mathtools}
\usepackage{tikz}
\usepackage[shortlabels]{enumitem}
\newtheorem{corollary}{Corollary}
\newtheorem{definition}{Definition}
\newtheorem{proposition}{Proposition}
\usepackage{cleveref}
\crefname{corollary}{Corollary}{Corollarys}
\crefname{definition}{Definition}{Definitions}
\crefname{proposition}{Proposition}{Propositions}
\newcommand{\ctauexc}[1]{\ensuremath{\mathsf{c}\text{-}\tau\text{-}\mathsf{exc} #1}}
\newcommand{\freegroup}[1]{\ensuremath{\mathsf{F}_{#1}}}
\newcommand{\mods}{\mathsf{mod}}
\title{Classifying Nakayama algebras with a braid group action on $\tau$-exceptional sequences}
\author{Maximilian Kaipel, Håvard U. Terland}
\date{Source: arXiv:2507.07608v2 (CC BY 4.0)}
\begin{document}
\maketitle

\noindent\emph{Source and licence.} Classifying Nakayama algebras with a braid group action on $\tau$-exceptional sequences. Version arXiv:2507.07608v2 (CC BY 4.0), distributed by arXiv under the Creative Commons Attribution 4.0 International licence, http://creativecommons.org/licenses/by/4.0/, as recorded in the arXiv licence metadata for this exact version and shown on the abstract page https://arxiv.org/abs/2507.07608v2. Full licence notice: \texttt{licenses/arxiv-2507.07608-CC-BY-4.0.txt}.\par\smallskip

\noindent\textit{Editorial scope.} Counted unit: the article's proposition prop:properties\_mutation\_braid\_group (source0.tex lines 444--454). The article's complete original proof of it is delivered with it (source0.tex lines 456--473, one \texttt{proof} environment of 18 lines). No mathematics is written, repaired or re-derived: the statement and the proof are the article's own lines, in the article's own notation, and no retained line is substituted or re-flowed.  Retained uncounted as the source-local context the proof needs: lines 274--276; lines 390--398; lines 400--402.  Nothing was omitted from any retained span, so the package carries no omission record.  No retained line contains a citation, so the wrapper carries no bibliography and no bibliography entry.  No retained line contains a figure environment, an image-inclusion command or drawing code, so no image is delivered and the package declares no asset dependency.  Editorial formatting only, in addition: an AMS \texttt{amsart} preamble that restates the article's own declarations and macros used by the retained lines, namely the environments \texttt{corollary}, \texttt{definition}, \texttt{proposition} on separate counters, together with the standard packages those lines need.  The retained environments print in this document as Corollary 1, Definition 1, Proposition 1 in order of appearance, whereas the article prints the same items with the numbering of its own full document.  The complete native source of the article as distributed in the arXiv e-print for this exact version is bundled unchanged as \texttt{sources/181357/source0.tex}.
\begin{corollary}\label{cor:tau_uniqueness_jasso_preserve}
   Let $(X,C)$ be a $\tau$-exceptional sequence. Then $X$ is uniquely determined by $C$ and $J(X,C)$.
\end{corollary}

\begin{definition}\label{def:braidgroup_relations}
    By the braid group relations on $n$ strands we mean the relations B1 and B2 below, where we assume $i,j \in \{1,2,\dots,n-1\}$: \begin{itemize}
         \item[\quad B1:] $\rho_i\rho_j = \rho_j\rho_i$ for $|i-j|\geq 2$.
        \item[\quad B2:] $\rho_{i}\rho_{i+1}\rho_{i} = \rho_{i+1}\rho_{i}\rho_{i+1}$ for $i\leq n-2$.
    \end{itemize}

    The braid group $B_n$ is defined as the group given by the generators $\{ \rho_1,\rho_2,\dots,\rho_{n-1} \}$ and the braid group relations B1 $\cup$ B2. 

\end{definition}

\begin{proposition}\label{prop:b1_relations_holds}
    The relations B1 are always respected by the canonical $\freegroup{n-1}$-action on $\ctauexc{\Lambda}$ in the sense that $\rho_i\rho_j\cdot A = \rho_j \rho_i \cdot A$ for any $A \in \ctauexc{\Lambda}$ and $i,j \in \{1,2,\dots,n-1\}$ where $|i-j|\geq 2$.
\end{proposition}

\begin{proposition}\label{prop:properties_mutation_braid_group}

    Consider the two criteria below for an algebra $\Lambda$.

    \begin{enumerate}
        \item[M1:] For any $\tau$-exceptional sequence $(B,C)$ we have that $\varphi(B,C) = (?,B)$.
        \item[M2:] For any indecomposable $\tau$-rigid module $X$ over $\Lambda$ and any $(B,C)$ which is $\tau$-exceptional both in $\mods \Lambda$ and in $J(X)$, we have $\varphi^{J(X)}(B,C) = \varphi(B,C)$.
    \end{enumerate}

    Let now $\mathcal{A}$ be a class of algebras closed under Jasso reductions such that criteria M1 and M2 above hold for any algebra in $\mathcal{A}$. Then for any algebra in this class, the mutation of complete $\tau$-exceptional sequences respects the braid group relations.
\end{proposition}

\begin{proof}
    Let $(A_1,A_2,\dots,A_n)$ be any complete $\tau$-exceptional sequence in $\ctauexc{\Lambda}$ for some $\Lambda$ in $\mathcal{A}$. We can assume that $t\geq 3$.

    Let $1 \leq i \leq t-2$ and consider the sequence $(A_i,A_{i+1},A_{i+2})$ in $J(A_{i+3},\dots,A_n)$. Since $\mathcal{A}$ is closed under Jasso reductions, we can consider $(A_i,A_{i+1},A_{i+2})$ to be a $\tau$-exceptional sequence in $\mods \Gamma$ for an algebra $\Gamma$ in $\mathcal{A}$. 

    We compute \begin{align}
        \varphi_1\circ \varphi_2 \circ \varphi_1(A_i,A_{i+1},A_{i+2}) &= \varphi_1 \circ \varphi_2(A^\star_{i+1},A_i,A_{i+2}) & \text{(By M1 in $J_\Gamma(A_{i+2})$)} \label{eq:line1_braid_relation}\\
        &= \varphi_1(A^\star_{i+1},A^\star_{i+2},A_i) & \text{(By M1 in $\mods \Gamma$)}\\
        &= (A^{\star\star}_{i+2},A^\star_{i+1},A_i) & \text{(By M1 in $J_\Gamma(A_{i})$)} \\ \notag \\
        \varphi_2\circ \varphi_1 \circ \varphi_2(A_i,A_{i+1},A_{i+2}) &= \varphi_2 \circ \varphi_1(A_{i},A^\dagger_{i+2},A_{i+1}) & \text{(By M1 in $\mods \Gamma$)}\\
        &= \varphi_2(A^{\dagger\dagger}_{i+2},A_{i},A_{i+1}) & \text{(By M1 in $J_\Gamma(A_{i+1})$)}\\
        &= (A^{\dagger\dagger}_{i+2},A^\star_{i+1},A_i) & \text{(By M2 and \cref{eq:line1_braid_relation})}
    \end{align} 


    
    Lastly by \cref{cor:tau_uniqueness_jasso_preserve} applied to $(A^{\dagger\dagger}_{i+2},A^\star_{i+1})$ and $(A^{\star\star}_{i+2},A^\star_{i+1})$ in $J(A_i)$ we see that $A^{\dagger\dagger}_{i+2} = A^{\star\star}_{i+2}$. We thus have that \[\varphi_i\circ \varphi_{i+1} \circ \varphi_i(A_1,A_2,\dots,A_n) = \varphi_{i+1} \circ \varphi_i \circ \varphi_{i+1}(A_1,A_2,\dots,A_n),\] demonstrating that mutation on $\ctauexc{\Lambda}$ respects the B2 relations in \cref{def:braidgroup_relations}. Since the mutation always respects the B1 relations by \cref{prop:b1_relations_holds}, the braid group relations on $n$ strands are satisfied.
\end{proof}

\end{document}
